\documentclass{exam}
\begin{document}

Nombre: \hrule
Ojo: no usar funciones precontruidas en haskell

\begin{enumerate}
\item Genere todas las permutaciones de [1,2,3,4,5,6] con un programa en haskell
\item Genere todas las conbinaciones de [1,2,3,4,5,6] con un programa en haskell. SUg usar subset
\item Genere el conjunto potencia de [1,2,3,4,5,6] (que representa un conjunto) en haskell. Sug usar subsets
\item Utice una dicision de casi tres partes de una lista para implementar
el algoritmo de fusion en haskell. Sug. inspiranse en el algorimo de fusion
\item Optimice el calculo de $3^n$ en haskell (tiempo logaritmico, no lineal).
\item Haga un oso de peluche en haskell Sug. Esto se realiza primero en contruir un texto que dibuje el oso y despues imprimirlo en pantalla a modo de texto.
\item Escriba el producto cartesiano de [1,2,3,4] y "abcd" (como lista de caracteres)
utilizando listas compactas.
\item Lee los datos de un archivo dato.txt: Un numero por linea, y almacenelos en una lista Sug. ver learnhaskellforyourgood
\item Dado los numeros de 1 al 10, que valores de x, y y z suman 15? (no importa
que se repitan, resuelva el problema usando listas compactas). Sug [(x,y,z) | x < ...,x+y+z==15]
\item utilize foldr para aplicar la funcion f(x,y)=x*y+y a la lista [1..10]
con semilla 1.
\item algoritmos de ordenamientos
\end{enumerate}
\end{document}
